\documentclass[11pt]{article}

\usepackage[margin=1.1in, top=1.0in, bottom=1.1in]{geometry}
\usepackage{graphicx}
\usepackage{caption}
\usepackage{subcaption}
\usepackage[colorlinks=true, linkcolor=blue, citecolor=blue, urlcolor=blue]{hyperref}
\usepackage{parskip}
\usepackage{setspace}
\usepackage{enumitem}
\usepackage{booktabs}
\usepackage{microtype}
\usepackage{textcomp}   % \textdegree
\usepackage{abstract}

\setstretch{1.15}

%% ─── Title block ─────────────────────────────────────────────────────────────
\title{%
  \textbf{BATPING: Weekly Progress Report}\\[4pt]
  \large Bidirectional Acoustics Traveltime Pulser for INdex Gradients\\[2pt]
  \normalsize Week of July 14--18, 2026
}
\author{Cyrus Urheim\\
        \small Department of Physics, Harvard University\\
        \small PRISE 2026
}
\date{July 20, 2026}

\begin{document}
\maketitle
\hrule
\vspace{4pt}

%% ─── Abstract ────────────────────────────────────────────────────────────────
\begin{abstract}
\noindent
This note summarises BATPING bench-prototype progress during the week of
July 14--18, 2026.
Three activities are reported: (1)~piezo-array transfer function
re-characterisation for both the UltraMic\,384k and ECM999 microphones;
(2)~new chirped-waveform recording sessions with both microphones and the
diagnosis and correction of three matched-filter pipeline failure modes that
produced jump discontinuities in the differential time-of-flight (TOF)
time series; and
(3)~a first attempt at a controlled temperature-change experiment using
a resistive dummy load on the optical rail, along with the issues identified
and the planned remediation.
\end{abstract}

\hrule
\bigskip

%% ─── Intro / context ─────────────────────────────────────────────────────────
\noindent
My main work this past week has been (1)~documenting and finalising all results
using the Dodotronic UltraMic (which digitises at 384\,kSa/s),
(2)~implementing the new ECM999 measurement microphone and running new
chirped-waveform recordings with it for repeatability, and (3)~attempting to
introduce a controlled temperature fluctuation along the acoustic path using
dummy-load resistors.

%% ─────────────────────────────────────────────────────────────────────────────
\section{Piezo Array Transfer Function Measurements}
%% ─────────────────────────────────────────────────────────────────────────────

I re-recorded transfer functions for both the piezo-array\,+\,UltraMic\,384k
and piezo-array\,+\,ECM999 systems, following your recommendation to sweep
frequency by frequency rather than using a broadband chirp.
For the UltraMic, each step used an 8\,ms burst of continuous sine (1.2\,V$_\mathrm{pp}$,
Savitzky--Golay smooth $w=71$), with a 20\,ms settle time between steps
(900 frequencies, 20--65\,kHz at 50\,Hz spacing).
For the ECM999, the SDG was set to continuous (BSWV) sine output at 0.5\,V$_\mathrm{pp}$
and the PicoScope captured a 50\,ms steady-state window per step with the
same settle time and frequency grid.
Both approaches yield per-step SNR of 30--60\,dB across the piezo passband.
Figure~\ref{fig:tf_comparison} compares the two resulting system transfer
functions.

\begin{figure}[h!]
  \centering
  \includegraphics[width=\textwidth]{tf_comparison_ultramic_ecm999.png}
  \caption{%
    Piezo-array system transfer function $H(f)$, normalised to the smoothed
    peak, for the UltraMic\,384k (\textit{left}; 8\,ms burst stepped-sine,
    1.2\,V$_\mathrm{pp}$) and the ECM999\,+\,SR560\,$\times$100 via
    PicoScope\,5444D (\textit{right}; BSWV continuous sine, 0.5\,V$_\mathrm{pp}$).
    Both sweeps used 900 steps from 20--65\,kHz at 50\,Hz spacing with 20\,ms
    settle time; responses are Savitzky--Golay smoothed ($w=71$).
    Thin traces show the raw per-step amplitude; bold traces show the smooth.
    Shaded bands and dashed lines indicate the $-3$\,dB and $-6$\,dB
    bandwidths relative to the smoothed peak.
    The solid blue vertical line marks the 40\,kHz BATPING carrier.}
  \label{fig:tf_comparison}
\end{figure}

Both systems are dominated by the piezo array's acoustic resonance, peaking
at 41.3--41.5\,kHz.
Both microphones yield very similar passbands: the UltraMic gives a
$-3$\,dB BW of 2.5\,kHz (40.4--42.9\,kHz) and $-6$\,dB BW of 4.3\,kHz,
while the ECM999 gives $-3$\,dB BW of 2.3\,kHz (40.3--42.6\,kHz) and
$-6$\,dB BW of 4.4\,kHz — essentially identical within measurement uncertainty.
The ECM999 shows a slightly broader passband, consistent with the difference
in measurement method (continuous sine vs.\ burst) and the SR560 preamplifier
shaping.
Notably, the UltraMic response is flat across the 20--65\,kHz sweep range
with no sharp rolloff, confirming its suitability for broadband characterisation,
while the ECM999 shows secondary features around 55--60\,kHz that are present
in both microphone measurements and therefore originate from the transmitting
piezo array rather than either microphone.

With the new transfer functions in hand, I recomputed the per-frequency
pre-compensation amplitude profiles for all chirped waveforms before
conducting the recording sessions described in Section~\ref{sec:sessions}.

%% ─────────────────────────────────────────────────────────────────────────────
\section{Chirped Waveform Sessions and Pipeline Improvements}
\label{sec:sessions}
%% ─────────────────────────────────────────────────────────────────────────────

Using the updated pre-compensation, I recorded two new sessions with each
microphone (UltraMic\,384k and ECM999).
Across all four sessions, sweeping burst durations of 2--4\,ms, bandwidths
of 6--20\,kHz, and pre-emphasis exponents $\alpha = 0.2$--$0.8$, we
continue to achieve per-ping LS-fit measurement uncertainty
$\delta(\Delta\tau)$ on the order of 2--5\,ns with both microphones.

The multi-session dataset also provided an opportunity to characterise
several failure modes that had been contributing to jump discontinuities
in the differential-TOF time series (Figure~\ref{fig:v1_failures}).

\begin{figure}[h!]
  \centering
  \includegraphics[width=0.88\textwidth]{recordings/20260716_015251_ecm999_full_sweep_v2_ecmHf/plots/bw_10dB_asym_2ms_comp_light/stubbs_pipeline_lsfit_v1/stubbs_lsfit.png}
  \caption{%
    Differential-TOF time series $\Delta\tau(t)$ from pipeline v1 for
    waveform \texttt{bw\_10dB\_asym\_2ms\_comp\_light}
    (ECM999\,+\,SR560, Jul.\,16 full-sweep session).
    The bottom panel shows multiple jump discontinuities scattered
    throughout the 30-second recording on an otherwise smooth underlying
    trend; the per-ping LS-fit uncertainty $\delta(\Delta\tau)$ during
    quiet stretches is $\sim$2\,ns, confirming that the discontinuities
    are artefactual rather than physical temperature fluctuations.}
  \label{fig:v1_failures}
\end{figure}

The jump discontinuities arose from three distinct failure modes, illustrated
in Figures~\ref{fig:diag_v3} and~\ref{fig:diag_v4}.
First, \textit{burst-window drift}: the waveform generator and digitiser
operate on independent clocks with a $\approx$187\,ppm frequency offset,
so over thousands of pings the acoustic burst gradually migrates within
its fixed analysis window; once the burst drifts far enough to be clipped
at the window boundary, the cross-correlation returns a spurious lag
(visible as periods labelled ``BAD'' in Figure~\ref{fig:diag_v3}, where
the reflected-signal trace in the analysis window is nearly zero).
Second, \textit{window overlap}: the outgoing and reflected burst analysis
windows were initially placed too close together, allowing the piezo
ring-down of the outgoing pulse to bleed into the reflected-pulse template.
Third, \textit{carrier-cycle slipping}: because the cross-correlation function
has sidelobes spaced one carrier period ($25\,\mu$s) apart, the estimator
occasionally locked onto the wrong sidelobe, producing a lag error that
is an integer multiple of $25\,\mu$s.
Each failure mode was individually diagnosed and patched in pipeline v2.

\begin{figure}[h!]
  \centering
  \includegraphics[width=\textwidth]{diag_a0.8_sm_4ms_v3.png}
  \caption{%
    Per-period diagnostic for waveform \texttt{a0.8\_sm\_4ms} under
    pipeline v1, showing five representative evaluation periods.
    \textit{Left}: drift-corrected full period with incident (pink) and
    reflected (blue) analysis windows.
    \textit{Centre}: co-added return template $T_R$ (red) overlaid on the
    actual received signal $m_\mathrm{ref}$ (grey) in the reflected window.
    \textit{Right}: Hilbert envelope of the reflected cross-correlation.
    In the three ``BAD'' periods (rows 3--5), accumulated clock drift has
    pushed the reflected burst outside the analysis window: the received
    signal in the centre panel is essentially zero, and the cross-correlation
    envelope is flat noise with no identifiable peak.}
  \label{fig:diag_v3}
\end{figure}

\begin{figure}[h!]
  \centering
  \includegraphics[width=\textwidth]{diag_a0.8_sm_4ms_v4.png}
  \caption{%
    Same diagnostic as Figure~\ref{fig:diag_v3} for the same five evaluation
    periods under the corrected pipeline v4.
    The adaptive period slicer now anchors each analysis window to the
    detected incident burst rather than using fixed time slots, so
    accumulated clock drift no longer displaces the reflected burst.
    All five periods are now ``GOOD'': the reflected signal fills the window,
    the cross-correlation has a sharp single peak, and the LS-fit uncertainty
    $\delta_R$ is 1.4--2.2\,ns across all periods including the formerly
    pathological evaluations 1309, 1312, and 1400.}
  \label{fig:diag_v4}
\end{figure}

With these corrections in place, pipeline v2 eliminates almost all
instances of measurement jitter across the 2\,m bench path
(Figure~\ref{fig:v2_clean}).

\begin{figure}[h!]
  \centering
  \includegraphics[width=0.88\textwidth]{recordings/20260716_015251_ecm999_full_sweep_v2_ecmHf/plots/bw_10dB_asym_2ms_comp_light/stubbs_pipeline_lsfit_v2/stubbs_lsfit.png}
  \caption{%
    Differential-TOF time series from pipeline v2 for the same waveform
    and session as Figure~\ref{fig:v1_failures}
    (\texttt{bw\_10dB\_asym\_2ms\_comp\_light}, ECM999).
    The jump discontinuities visible in v1 are eliminated; the time series
    follows a smooth physical trend with the per-ping uncertainty
    $\delta(\Delta\tau)$ reduced to $\sim$2\,ns.}
  \label{fig:v2_clean}
\end{figure}

%% ─────────────────────────────────────────────────────────────────────────────
\section{Controlled Temperature Experiment: Dummy-Load Resistor}
%% ─────────────────────────────────────────────────────────────────────────────

On Thursday and Friday I ran a first controlled temperature-change test using
a 25\,W RF dummy-load resistor (XRDS-RF, 50\,$\Omega$, BNC connector)
mounted at the midpoint of the 1\,m Thorlabs optical breadboard, directly
beneath the acoustic path.
DC power was supplied by a Keithley 2231A-30-3 triple-channel supply found
in the lab, and BATPING was recorded continuously using
\texttt{picoscope\_batping\_record.py} while voltage was slowly ramped.
Two power levels were explored: 10\,V ($\approx$2\,W, 8\% of the 25\,W
rating) and 30\,V (18\,W, 72\% of rating, the Keithley's maximum output).
The resistor did not heat up appreciably at 10\,V after 15\,minutes;
it was warm to the touch at 30\,V.
In both cases the recordings showed nanosecond-scale uncertainty during
quiet stretches, but with frequent large jump discontinuities.

Two issues were identified post-hoc.
First, \textit{acoustic scintillation}: at 18\,W the resistor drove a strong
convective plume above the rail.
Turbulent air scattered the acoustic wavefront --- the acoustic analogue of
optical dome seeing --- causing the received signal amplitude to fluctuate
by more than an order of magnitude between 30-second analysis windows.
Windows with near-zero amplitude yield meaningless lag estimates regardless
of pipeline quality.
Second, \textit{incorrect PicoScope sample rate}: the TIMEBASE register had
been set using the formula for 8-bit ADC resolution, whereas the unit was
operating in 12-bit mode; this resulted in an actual sample rate of
251\,kSa/s (6.3 samples per 40\,kHz carrier cycle) rather than the intended
1\,MSa/s (25.6 samples/cycle).

Both issues have been corrected.
The recording script now uses the proper 12-bit TIMEBASE value (verified by
direct SDK query to the connected scope), and the next heat test will use
a much lower power level of 5--10\,V (0.5--2\,W).
At these levels the thermal gradient should develop slowly and without a
disruptive convective plume; at 1\,W, an air-column temperature rise of
$\sim$0.1--0.5\,\textdegree C over 5--10\,minutes is expected --- well
within BATPING's $\sim$0.05\,\textdegree C per-ping sensitivity.
I plan to re-record at multiple power levels tomorrow and will send an
update then.

\bigskip
\hrule
\vspace{4pt}
\noindent\small\textit{%
  All analysis code, recording scripts, and figures are version-controlled
  in the project repository.
  Report prepared \today.}

\end{document}
